\documentclass[11pt]{article}
\usepackage[a4paper,margin=0.96in]{geometry}
\usepackage{amsmath,amssymb,amsthm,mathtools}
\usepackage{booktabs,array,graphicx,xcolor,xurl}
\usepackage[hidelinks]{hyperref}
\hypersetup{pdftitle={The Rarest Target Sets Passive Memory},pdfauthor={Lluis Eriksson},pdfsubject={Exact McMillan degree of spectral routing words}}
\usepackage[nameinlink,capitalise]{cleveref}
\crefname{lemma}{Lemma}{Lemmas}
\Crefname{lemma}{Lemma}{Lemmas}
\crefname{proposition}{Proposition}{Propositions}
\Crefname{proposition}{Proposition}{Propositions}
\crefname{corollary}{Corollary}{Corollaries}
\Crefname{corollary}{Corollary}{Corollaries}
\usepackage{microtype,etoolbox}
\AtBeginEnvironment{thebibliography}{\fontsize{8}{9}\selectfont\setlength{\itemsep}{-2pt}\setlength{\parsep}{0pt}}

\newtheorem{theorem}{Theorem}[section]
\newtheorem{lemma}[theorem]{Lemma}
\newtheorem{corollary}[theorem]{Corollary}
\newtheorem{proposition}[theorem]{Proposition}
\theoremstyle{definition}\newtheorem{definition}[theorem]{Definition}\newtheorem{example}[theorem]{Example}
\theoremstyle{remark}\newtheorem{remark}[theorem]{Remark}\newtheorem{limitation}[theorem]{Limitation}
\newcommand{\C}{\mathbb C}\newcommand{\D}{\mathbb D}\newcommand{\T}{\mathbb T}
\newcommand{\ran}{\operatorname{ran}}\newcommand{\mcdeg}{\deg_{\mathrm{McM}}}
\newcommand{\norm}[1]{\lVert#1\rVert}\newcommand{\A}{\mathcal A}

\title{The Rarest Target Sets Passive Memory:\\An Exact Routing-Word Degree Law}
\author{Lluis Eriksson}
\date{11 August 2026}

\begin{document}
\maketitle

\begin{abstract}
We determine the exact passive memory required to route one fixed input
subspace repeatedly among orthogonal spectral destinations.  Let a routing
word $w$ of length $L$ use target $a$ exactly $n_a$ times, and let every
target and the input have dimension $k$.  Among all square finite
rational-inner transfer matrices, regular at the prescribed distinct
boundary nodes and realizing the subspace routing constraints, we prove
\[
                  d_{\min}=k\bigl(L-\min_{a\in\A(w)}n_a\bigr).
\]
The lower bound is a zero budget: a compressed $k\times k$ minor vanishes to
order at least $k$ whenever the route points elsewhere, while its total zero
budget cannot exceed McMillan degree.  The matching construction is
explicit.  Node polynomials are normalized by scalar Fej\'er--Riesz
factorization to form a rational-inner routing column; a degree-preserving
lossless completion and tensor amplification produce a square network of
the asserted degree.

Thus passive memory is governed by the rarest destination and, within this
model, is invariant under every permutation of the routing word.  Switch
count need not even order the cost: for scalar channels, the cyclic word
\texttt{AAAB} has two switches and degree three, whereas \texttt{ABAB} has
four switches and degree two.  The law recovers constant routing, one change
along a linearly ordered node sequence, and fully distinct fan-out as exact
limiting cases.  It also yields
the optimal integrated Wigner--Smith delay.  A deterministic compiler and
an independent verifier replay 1,024 random words, all 1,155 word classes
through length seven modulo alphabet renaming, permutation tests, and
clustered-node stress fixtures.  The machine-readable certificate and the
paper are linked to an immutable public artifact.
\end{abstract}

\noindent\textbf{Keywords:} rational inner matrix; McMillan degree;
spectral routing; Fej\'er--Riesz factorization; paraunitary completion;
passive memory; Wigner--Smith delay; routing word.

\section{A word, not a path length}

A lossless finite-dimensional discrete-time network is represented by a
square rational matrix $S$ analytic on $\D$ and unitary almost everywhere on
$\T$.  Its McMillan degree is the dimension of a minimal conservative
realization.  Fix distinct nodes $\zeta_1,\ldots,\zeta_L\in\T$, a
$k$-plane $X\subset\C^N$, and mutually orthogonal $k$-planes
$Y_a\subset\C^N$ indexed by a finite alphabet $\A$.  A word
$w=w_1\cdots w_L\in\A^L$ prescribes
\begin{equation}\label{eq:task}
                         S(\zeta_i)X=Y_{w_i},\qquad 1\le i\le L.
\end{equation}
Only the subspaces are prescribed; bases and phases are free.  We minimize
$\mcdeg S$ over square finite rational-inner $S$ regular at every node.
Symbols not occurring in $w$ are discarded, and
\[
 n_a:=\#\{i:w_i=a\},\qquad n_*:=\min_{a\in\A(w)}n_a.
\]

Previous exact results isolated the cost of one switch and of fully
distinct orthogonal fan-out \cite{ErikssonSwitch2026,ErikssonGlobalFanout2026};
the direct-sum sequel showed that maximal fan-out is geometrically generic
\cite{ErikssonDirectSum2026}.  Repeated targets create a different question:
does memory count transitions, distinct targets, or interpolation
constraints?  The answer is a sharp occupancy statistic.

\begin{theorem}[Rarest-target law]\label{thm:main}
Under the hypotheses above,
\begin{equation}\label{eq:law}
             d_{\min}(w;k)=k(L-n_*).
\end{equation}
The value is independent of the locations of the distinct nodes and of the
order of the letters.
\end{theorem}

The order-independence in \cref{thm:main} concerns minimum state dimension,
not every dynamical observable of a particular realization.  It is already
strong enough to separate memory from switch count.

\begin{table}[b]
\centering\small
\caption{Scalar examples.  Cyclic switches count changes across all four
adjacencies, including the last-to-first edge.}
\label{tab:words}
\begin{tabular}{@{}lccccl@{}}\toprule
Word & $L$ & occupancies & cyclic switches & $d_{\min}$ & interpretation\\\midrule
\texttt{AAAA} &4&(4)&0&0&constant route\\
\texttt{AAAB} &4&(3,1)&2&3&rare excursion\\
\texttt{AABB} &4&(2,2)&2&2&two contiguous blocks\\
\texttt{ABAB} &4&(2,2)&4&2&alternating route\\
\texttt{ABCD} &4&(1,1,1,1)&4&3&full fan-out\\
\texttt{AABBCC}&6&(2,2,2)&3&4&balanced three-target route\\
\bottomrule\end{tabular}
\end{table}

\begin{figure}[t]
\centering
\includegraphics[width=0.98\linewidth]{figures/routing_word_memory.pdf}
\caption{Exact scalar memory at length four.  The law depends on the rarest
occupancy, not monotonically on cyclic switch count.  In particular,
\texttt{AAAB} costs more than \texttt{ABAB}.}
\label{fig:separation}
\end{figure}

\section{A zero budget for compressed minors}

Choose isometric frames $X_0$ for $X$ and $U_a$ for $Y_a$.  For an
interpolant $S$, define the scalar compressed minor
\begin{equation}\label{eq:minor}
                         g_a(z)=\det(U_a^*S(z)X_0).
\end{equation}
At a node carrying $a$, the matrix inside the determinant is unitary, up to
the freely chosen input and output gauges; hence $g_a(\zeta_i)\ne0$.  At a
node carrying $b\ne a$, orthogonality gives
$U_a^*S(\zeta_i)X_0=0$.  All $k^2$ entries therefore vanish there, and the
determinant has a zero of order at least $k$.

We record the degree fact that converts these forced zeros into states.

\begin{lemma}[Minor zero budget]\label{lem:budget}
Let $S$ be a square finite rational-inner matrix of McMillan degree $d$, and
let $U,V$ be constant isometries with $k$ columns.  If
$g(z)=\det(U^*S(z)V)$ is not identically zero, then the total multiplicity of
its boundary zeros at regular points of $S$ is at most $d$.
\end{lemma}

\begin{proof}
Factor $S$ into a constant unitary and $d$ elementary rank-one
Blaschke--Potapov factors, counting multiplicity \cite{Potapov1955}.  Write
one factor as $E(z)=I+(b(z)-1)P$, where $P$ is rank one.  Its exterior-power
action is explicitly
\begin{equation}\label{eq:wedge-factor}
 \bigwedge\nolimits^k E(z)=I+(b(z)-1)\Pi_P,
\end{equation}
where $\Pi_P$ projects onto the wedge subspace whose elementary wedges
contain the direction $\ran P$.  Thus each factor is affine in its scalar
Blaschke coordinate.  On $\bigwedge^k\C^N$, $g$ is the matrix coefficient
$\langle\bigwedge^kU,(\bigwedge^kS)\bigwedge^kV\rangle$ and is therefore
multiaffine in the $d$ scalar coordinates.  After clearing their $d$
denominators, every matrix coefficient has a numerator of degree at most
$d$.  A nonzero such numerator has at most $d$ zeros, with boundary
multiplicity included.  Constant unitary factors and cancellations can only
decrease this number.
\end{proof}

The exterior-power formulation is useful because it is insensitive to
target-frame gauges.  Replacing $U_a$ by $U_aR_a$ and $X_0$ by $X_0R_0$
multiplies $g_a$ by the nonzero constant
$\overline{\det R_a}\det R_0$; it cannot create or remove zeros.  Nor can a
pole be hidden at an interpolation node: regularity is explicitly part of
the comparison class.  These two details prevent a calibrated-frame or
removable-singularity loophole in the lower bound.

\begin{proposition}[Lower bound]\label{prop:lower}
Every interpolant of \eqref{eq:task} has degree at least $k(L-n_*)$.
\end{proposition}
\begin{proof}
For a fixed used symbol $a$, the minor \eqref{eq:minor} is nonzero at its
$n_a$ own nodes and has a zero of order at least $k$ at each of the other
$L-n_a$ nodes.  By Lemma~\ref{lem:budget},
$\mcdeg S\ge k(L-n_a)$.  Taking an $a$ of minimum occupancy proves the
claim.
\end{proof}

This proof also explains why counting adjacent changes loses information.
The obstruction is global: every destination monitors all nodes at which it
is absent.

\section{The optimal compiler}

The lower bound is attained by a construction that first solves the scalar
alphabet problem.  Put $m=|\A(w)|$ and
\begin{equation}\label{eq:pa}
                 p_a(z)=\prod_{i:w_i\ne a}(z-\zeta_i),
                 \qquad \deg p_a=L-n_a\le d:=L-n_*.
\end{equation}
The Laurent polynomial
\begin{equation}\label{eq:R}
                           R(z)=\sum_{a\in\A(w)}|p_a(z)|^2,
                           \qquad z\in\T,
\end{equation}
is strictly positive.  Indeed, away from the interpolation nodes none of the
$p_a$ vanishes; at $\zeta_i$, precisely the polynomial indexed by $w_i$ is
nonzero.  Scalar Fej\'er--Riesz factorization gives an outer polynomial $h$
of degree at most $d$, zero-free in $\D$, such that
$|h|^2=R$ on $\T$ \cite{DritschelRovnyak2009}.  In particular,
$h(\zeta_i)\ne0$ at every interpolation node, so the quotient below is
regular at all prescribed points.  Therefore
\begin{equation}\label{eq:f}
                       f(z)=\frac{1}{h(z)}(p_a(z))_{a\in\A(w)}
\end{equation}
is an $m\times1$ rational-inner column.  At $\zeta_i$ all its coordinates
except $w_i$ vanish and the surviving coordinate has modulus one.

Rectangular paraunitary functions admit square lossless completions with the
same realization state dimension \cite[Theorem 2.1(ii)]{AlpayJorgensenLewkowicz2016};
related polynomial spectral completions appear in
\cite{EphremidzeLagvilava2011}.  The cited theorem uses the exterior-disk
stability convention.  We apply it to the paraconjugate column
$f^\sharp(z)=f(1/\overline z)^*$ and paraconjugate the completed matrix back;
this returns the disk-analytic rational-inner convention used here without
changing state dimension.  Complete $f$ in this way to an $m\times m$
rational-inner $V$ without increasing its degree.  The scalar case of
Proposition~\ref{prop:lower} forces $\mcdeg V=d$, so no concealed cancellation is
possible.

For rank $k$, take $V\otimes I_k$.  Its McMillan degree is $kd$: each scalar
elementary factor is repeated in $k$ independent channels.  Constant
unitaries align the first input $k$-plane and the $m$ coordinate target
$k$-planes with $X$ and $Y_a$; an identity block embeds the construction if
$N>mk$.  Constant operations change no degree.  We have proved the matching
upper bound and hence \cref{thm:main}.

For completeness, the compiler can be stated as finite exact algebra before
any numerical realization step:
\begin{enumerate}
\item delete unused symbols and count $n_a$;
\item form the node polynomials \eqref{eq:pa};
\item compute the positive trigonometric polynomial $R$ in \eqref{eq:R};
\item choose its outer Fej\'er--Riesz factor $h$ and form $f$ in
      \eqref{eq:f};
\item apply a degree-preserving square lossless completion to $f$;
\item tensor by $I_k$ and conjugate by constant input/output unitaries.
\end{enumerate}
The first four steps already give a minimal scalar routing column.  The fifth
does not add states: in realization language it appends input and output
channels around the same state matrix.  The sixth repeats the scalar state
space $k$ times.  This makes the upper proof constructive at the transfer,
realization, and network-size levels.

\begin{remark}[Algorithmic content]
The proof is a compiler: construct the $p_a$, spectrally factor their summed
energy, form $f$, complete losslessly, tensor, and align.  No nonlinear
search over network parameters is required.
\end{remark}

\section{Consequences and separations}

\begin{corollary}[Exact special cases]\label{cor:special}
The law gives: (i) constant routing, $d_{\min}=0$; (ii) two distinct nodes
routed to two targets, $d_{\min}=k$; (iii) $L$ distinct targets,
$d_{\min}=k(L-1)$; and (iv) balanced use of $m$ targets when $m\mid L$,
$d_{\min}=kL(1-1/m)$.
\end{corollary}

\begin{example}[More switches, less memory]\label{ex:switch}
For $k=1$, \texttt{AAAB} has two cyclic switches, occupancies $(3,1)$, and
degree $3$.  The alternating word \texttt{ABAB} has four cyclic switches,
occupancies $(2,2)$, and degree $2$.  Thus switch count does not monotonically
predict passive memory.  Both examples use the same two target lines, so
target span alone also fails to determine the answer.
\end{example}

\begin{proposition}[Complete combinatorial quotient]\label{prop:quotient}
For fixed $L$ and $k$, two routing words with orthogonal targets have the
same minimum degree whenever their least nonzero occupancies agree.  Every
integer in
\[
                 k\lceil L/2\rceil,\ k(\lceil L/2\rceil+1),
                 \ldots,\ k(L-1)
\]
occurs among words using at least two targets, while the one-target word has
degree zero.
\end{proposition}
\begin{proof}
The first statement is \cref{thm:main}.  With two used letters, the smaller
occupancy can be any $r\in\{1,\ldots,\lfloor L/2\rfloor\}$, giving degree
$k(L-r)$.  These are precisely the displayed values.  More letters may
repeat values but cannot fill a missing value because the degree is still
$k(L-n_*)$.
\end{proof}

The discontinuity between one and two symbols is genuine.  A word may use a
second target only once; that single rare visit changes the minimum from
zero to $k(L-1)$.  Exact routing asks for a full $k$-plane zero at every
other node, so rarity is expensive rather than cheap.

The theorem may be read as a law for a labelled walk on any graph whose
vertices are represented by mutually orthogonal output planes: once the
word of visited vertices is fixed, the minimum state dimension depends only
on the least visit count.  This is not a claim for nonorthogonal graph
embeddings or edge-constrained internal dynamics.

For a differentiable unitary boundary value, use the Wigner--Smith matrix
\begin{equation}\label{eq:Q}
                   Q(\theta)=-iS(e^{i\theta})^*\frac{d}{d\theta}S(e^{i\theta}).
\end{equation}
The argument principle for a finite rational-inner matrix yields
$\int_0^{2\pi}\operatorname{tr}Q(\theta)d\theta=2\pi\mcdeg S$.  Hence:

\begin{corollary}[Optimal integrated delay]\label{cor:delay}
Among networks realizing the routing word,
\begin{equation}\label{eq:delay}
 \min_S\int_0^{2\pi}\operatorname{tr}Q(\theta)d\theta
       =2\pi k(L-n_*).
\end{equation}
\end{corollary}

The equality is topological and global.  It does not prescribe a unique
pointwise delay profile, peak delay, or physical dwell-time distribution.

\subsection{Conservative-state interpretation}

If $S(z)=D+zC(I-zA)^{-1}B$ is a minimal conservative realization, then
$\mcdeg S=\dim A$.  The theorem therefore counts internal lossless state
coordinates, not merely polynomial coefficients.  In this representation
the lower proof says that the same state space must supply $k(L-n_a)$
independent boundary-zero multiplicities for every target monitor $a$.
The upper construction supplies exactly the largest of these demands and
shares those states among the other monitors.  This max-sharing mechanism is
why the answer is
\[
                       \max_a k(L-n_a)=k(L-n_*)
\]
rather than the sum over targets.

The identity also separates architecture from calibration.  Constant
unitary pre- and post-networks can change the physical port basis at zero
McMillan cost; only the frequency-dependent inner core contributes states.
Conversely, allowing active gain or nonlossless cancellation would change
the class and invalidate both the delay identity and the Potapov zero
budget.

\section{Executable certificate}

The repository contains two deliberately separate programs.  The producer
constructs $p_a$, performs scalar spectral factorization by reciprocal-root
selection, evaluates the inner column on a dense boundary grid, checks every
node, freezes a JSON certificate, and creates \cref{fig:separation}.  The
verifier imports no producer code: it validates the schema and digest,
recomputes the combinatorial law, and enforces numerical tolerances.

The frozen campaign contains 1,024 seeded random words of lengths two
through nine; every one of the 1,155 restricted-growth word classes through
length seven (alphabet renaming removed); 256 permutation checks; ten
curated cases; and four \texttt{ABAB} fixtures with a closest prescribed
angular gap decreasing from $10^{-1}$ to $10^{-4}$.  All pass.  Across the
full campaign the largest boundary column-isometry residual is
$1.63\times10^{-5}$, the largest wrong-target leakage is
$8.0\times10^{-11}$, and the smallest desired amplitude is $0.999996$.
For the exhaustive exact-grid subcampaign the corresponding isometry and
leakage maxima are $1.1\times10^{-13}$ and $4.4\times10^{-15}$.

The larger residual in clustered random fixtures is a conditioning warning
for root-based floating-point factorization, not an exception to the exact
identity.  The certificate fails closed at stated tolerances; it is evidence
that the constructive formulas were transcribed and replayed correctly, not
a substitute for the proof.

The producer additionally stores the word, node angles, occupancies,
predicted degree, spectral-factor residual, boundary isometry residual,
desired amplitude, and wrong-coordinate leakage for each curated and stress
case.  A SHA-256 digest covers the canonical JSON payload before the final
pass flag is appended.  The independent verifier recomputes that digest and
the degree formula from the stored words.  This division catches both
floating-point regression in the compiler and accidental mutation of the
frozen evidence.

\begin{table}[t]
\centering\small
\caption{Theorem-to-artifact map.  Paths are relative to the public repository.}
\label{tab:artifacts}
\begin{tabular}{@{}p{0.30\linewidth}p{0.62\linewidth}@{}}\toprule
Claim & Executable witness\\\midrule
Occupancy law and compiler & \path{research/routing_word_memory_certificate.py}\\
Frozen campaign and residuals & \path{results/routing_word_memory/certificate.json}\\
Independent replay & \path{verification/verify_routing_word_memory.py}\\
Release-wide integrity gate & \path{verification/run_routing_word_memory_release_checks.py}\\
Switch-count separation & producer-generated \path{paper_routing_word_memory/figures/routing_word_memory.pdf}\\
\bottomrule\end{tabular}
\end{table}

\section{Scope, conditioning, and relation to prior work}

Boundary matrix interpolation and McMillan-degree minimization have a deep
general theory \cite{AntoulasEtAl1990,BallKim1994}, and modern subspace
all-pass design can prescribe unitary data and group-delay blocks
\cite{BharathEtAl2023}.  Fej\'er--Riesz factorization and lossless completion
used here are classical.  The contribution is their combination with the
minor zero budget to expose a closed occupancy law for repeated orthogonal
subspace routing.  A targeted search of these literatures did not locate the
formula \eqref{eq:law}; we therefore state novelty for this precise theorem,
not for rational-inner interpolation in general.

\begin{limitation}[Exact hypotheses]
The theorem assumes distinct boundary nodes, exact subspace interpolation,
mutually orthogonal equal-rank targets, and no prescribed frame inside a
target.  Repeated nodes require jet data.  Nonorthogonal targets remove the
forced-zero mechanism.  Fully calibrated frames constrain phases discarded
by \eqref{eq:task}.  Active, nonrational, or nonlossless devices lie outside
the comparison class.
\end{limitation}

\subsection{Why the hypotheses are structural}

Distinct nodes ensure that the polynomials \eqref{eq:pa} encode independent
point constraints.  If two nodes collide, value interpolation becomes
inconsistent unless their labels agree; a meaningful limit with changing
labels requires derivatives and hence a jet problem.  Orthogonality is what
makes the wrong-target compression identically zero.  For merely separated
targets it is small but nonzero, so the discrete zero count disappears even
though approximate lower bounds may survive.  Equal target ranks allow one
tensor amplification; heterogeneous ranks would attach different
determinant multiplicities to different absences and require a new packing
argument.

Free target gauges are equally important.  At a correct node the construction
delivers the requested plane with a node-dependent phase.  Fixing a complete
unitary frame would turn those phases into data and can increase degree.
The present theorem is consequently a statement about routing information,
not waveform calibration.

The three fan-out papers are complementary rather than duplicate statements:
the orthogonal paper supplies the distinct-target extremum, the direct-sum
paper identifies its generic geometric persistence and collision strata,
and the present word law classifies repeated orthogonal targets.  A unified
routing-table theory combining these strata is a natural future synthesis,
but none of the three theorems alone subsumes the other two.

Node positions do not alter exact degree, but they affect numerical
conditioning and realization coefficients.  Likewise, an arbitrarily small
generic perturbation of a repeated target can turn repetitions into distinct
direct-sum constraints and raise the exact degree, as the collision analysis
in the direct-sum paper shows \cite{ErikssonDirectSum2026}.  Robust
approximate routing is therefore a different optimization problem.

Three extensions are natural.  First, determine a stability law with
explicit angular and subspace-error margins.  Second, replace orthogonality
by a graph of prescribed intersections and identify which zero budgets
survive.  Third, characterize local or peak delay among the degree-optimal
completions.  None is needed for the exact word law proved here.

\section{Conclusion}

Repeated spectral routing has a complete state-counting law.  Each target
asks the network to vanish, in a $k$-fold determinant sense, at every node
assigned elsewhere.  The least visited target therefore imposes the largest
zero budget.  Fej\'er--Riesz normalization realizes that budget without
waste.  The resulting identity $d_{\min}=k(L-n_*)$ is simultaneously a
minimal-degree theorem, an explicit compiler, and a global delay law.  Its
most visible consequence is also the simplest: a route with more switches
can require less memory.

\paragraph{Reproducibility.}
The canonical paper, source, frozen JSON certificate, independent verifier,
and CI release gate are archived at
\href{https://github.com/lluiseriksson/finite-sample-spectral-certificates/releases/tag/v2.5-routing-word-memory-final}{the immutable public release}.

\bibliographystyle{plain}
\bibliography{references}
\end{document}
